\documentclass[a4paper,12pt]{article}

\usepackage[T2A]{fontenc}
\usepackage[utf8]{inputenc}
\usepackage[russian]{babel}
\usepackage{amsmath,amssymb,mathtools}
\usepackage[a4paper,margin=2cm,headheight=15pt,footskip=24pt]{geometry}
\usepackage{microtype}
\usepackage{array,booktabs,tabularx}
\usepackage{xcolor}
\usepackage{hyperref}
\usepackage{tikz}
\usepackage{fancyhdr}
\usepackage{titlesec}

\definecolor{Accent}{RGB}{48,78,110}
\definecolor{Muted}{RGB}{92,101,110}
\definecolor{Result}{RGB}{0,112,110}

\hypersetup{
  colorlinks=true,
  linkcolor=Accent,
  urlcolor=Accent,
  pdfauthor={Лёвин Артём Александрович},
  pdftitle={Разбор домашней работы по математике}
}

\setlength{\parindent}{0pt}
\setlength{\parskip}{0.4em}
\setlength{\abovedisplayskip}{7pt}
\setlength{\belowdisplayskip}{7pt}
\setlength{\abovedisplayshortskip}{5pt}
\setlength{\belowdisplayshortskip}{5pt}
\renewcommand{\arraystretch}{1.22}
\titlespacing*{\subsection}{0pt}{1.4ex plus .2ex minus .1ex}{0.45ex}
\setlength{\emergencystretch}{2em}

\pagestyle{fancy}
\fancyhf{}
\fancyfoot[L]{\footnotesize\color{Muted}Лёвин Артём Александрович эксклюзивно для Ксении Васильченко}
\fancyfoot[R]{\footnotesize\color{Muted}\thepage}
\renewcommand{\headrulewidth}{0pt}
\renewcommand{\footrulewidth}{0pt}

\newcolumntype{Y}{>{\raggedright\arraybackslash}X}
\newcommand{\ErrorLabel}{\textcolor{red!75!black}{\textbf{Ошибка:}}}
\newcommand{\KeyIdea}{\textcolor{Accent}{\textbf{Ключевая идея:}}}
\newcommand{\Outcome}{\textcolor{Result}{\textbf{Итог:}}}

\begin{document}

\begin{titlepage}
\thispagestyle{empty}
\centering
\vspace*{0.16\textheight}
{\Huge\bfseries Разбор домашней работы\par}
\vspace{0.7cm}
{\Large Математика: НОД и уравнения\par}
\vspace{1.6cm}
\begin{tabular}{@{}rl@{}}
\textbf{Ученица:} & Ксения Васильченко \\
\textbf{Преподаватель:} & Лёвин Артём Александрович \\
\textbf{Дата:} & 23 сентября 2026 года
\end{tabular}
\vfill
{\small\color{Muted}Лёвин Артём Александрович эксклюзивно для Ксении Васильченко\par}
\vspace{0.55cm}
\begin{tikzpicture}[x=1cm,y=1cm]
  \draw[Accent!70,line width=0.8pt]
  (0,0) .. controls (1.5,0.45) and (3.1,-0.45) .. (4.7,0)
        .. controls (6.1,0.4) and (7.5,-0.22) .. (8.8,0.08);
\end{tikzpicture}
\vspace{0.25cm}
\end{titlepage}

\newpage
\section*{Сводка по заданиям, требующим доработки}
\small
\begin{tabularx}{\linewidth}{@{}p{0.10\linewidth}p{0.18\linewidth}Yp{0.18\linewidth}p{0.18\linewidth}@{}}
\toprule
\textbf{№} & \textbf{Тема} & \textbf{Суть недочёта} & \textbf{Ваш результат} & \textbf{Требуемый результат} \\
\midrule
\hyperref[task:2]{2} & НОД по разложениям & Числа разложены правильно, но НОД для $72$ и $120$ не найден: далее запись переходит к числам из следующего задания. & НОД не указан & $24$ \\
\addlinespace
\hyperref[task:5a]{5а} & Линейное уравнение & Корень найден верно, но отсутствует обязательная проверка подстановкой. & $x=20$ & $x=20$ и проверка \\
\addlinespace
\hyperref[task:5b]{5б} & Переменная в знаменателе & Не записано ограничение $x\ne2$ до решения; после нахождения корня отсутствует обязательная проверка. & $x=6$ & $x=6$, $x\ne2$, и проверка \\
\bottomrule
\end{tabularx}
\normalsize

\newpage
\thispagestyle{empty}
\vspace*{\fill}
\begin{center}
{\Huge\bfseries Разбор ошибок}
\end{center}
\vspace*{\fill}

\newpage
\phantomsection\label{task:2}
\section*{Задание 2}

\subsection*{Текст задания}
Разложите числа $72$ и $120$ на простые множители. Найдите
\[
\text{НОД}(72,120).
\]

\subsection*{Ваше решение}
В работе оба исходных числа разложены на простые множители. По записанным шагам получается
\[
72=2\cdot2\cdot2\cdot3\cdot3=2^3\cdot3^2,
\]
\[
120=2\cdot2\cdot2\cdot3\cdot5=2^3\cdot3\cdot5.
\]
После этого вместо вычисления НОД чисел $72$ и $120$ начинается запись с числами $96$ и $144$, относящимися уже к следующему заданию.

\subsection*{Ваш ответ}
Значение $\text{НОД}(72,120)$ не указано.

\subsection*{Ошибка}
\ErrorLabel\ решение не доведено до требуемого результата: после верных разложений не найден НОД исходных чисел.

\subsection*{Где именно возникает ошибка}
Последний полностью правильный шаг ---
\[
72=2^3\cdot3^2,
\qquad
120=2^3\cdot3\cdot5.
\]
На этом этапе уже есть все данные для ответа. Следующим шагом нужно выбрать простые множители, которые входят в оба разложения, и для каждого общего множителя взять меньшую степень.

Вместо этого решение переходит к числам $96$ и $144$. Эти вычисления могут быть полезны для следующей задачи, но не отвечают на вопрос о НОД чисел $72$ и $120$.

\subsection*{Почему это неверно}
Разложение на простые множители --- только подготовительный этап. Чтобы получить НОД, нужно выделить общую часть двух разложений.

Для двойки в обоих числах есть степень $3$, поэтому в НОД входит $2^3$. Для тройки степени равны $2$ и $1$, поэтому берём меньшую степень $1$. Пятёрка есть только в разложении числа $120$, поэтому в НОД не входит.

Следовательно,
\[
\text{НОД}(72,120)=2^3\cdot3=8\cdot3=24.
\]

\subsection*{Ключевая идея}
\KeyIdea\ НОД по разложениям составляют только из общих простых множителей; для каждого из них берут меньшую степень.

\subsection*{Исправление от последнего правильного шага}
Сохраняем уже полученные разложения:
\[
72=2^3\cdot3^2,
\qquad
120=2^3\cdot3\cdot5.
\]
Общая часть равна
\[
2^3\cdot3.
\]
Поэтому
\[
\text{НОД}(72,120)=24.
\]

\subsection*{Полное правильное решение}
Разложим число $72$ на простые множители:
\[
72:2=36,
\]
\[
36:2=18,
\]
\[
18:2=9,
\]
\[
9:3=3,
\]
\[
3:3=1.
\]
Значит,
\[
72=2^3\cdot3^2.
\]

Разложим число $120$:
\[
120:2=60,
\]
\[
60:2=30,
\]
\[
30:2=15,
\]
\[
15:3=5,
\]
\[
5:5=1.
\]
Значит,
\[
120=2^3\cdot3\cdot5.
\]

Берём общие множители с меньшими степенями:
\[
\text{НОД}(72,120)=2^3\cdot3=24.
\]

\subsection*{Проверка}
Число $24$ делит оба исходных числа:
\[
72:24=3,
\qquad
120:24=5.
\]
После деления остаются числа $3$ и $5$, у которых нет общего делителя больше единицы. Значит общий делитель увеличить уже нельзя, и НОД действительно равен $24$.

\subsection*{Итог}
\Outcome\ разложения чисел выполнены правильно. Не хватило последнего обязательного шага: из общей части разложений нужно получить $\text{НОД}(72,120)=24$.

\subsection*{Задача на закрепление}
Разложите числа $84$ и $126$ на простые множители и найдите их наибольший общий делитель.

\newpage
\phantomsection\label{task:5a}
\section*{Задание 5а}

\subsection*{Текст задания}
Решите уравнение. Выполните проверку:
\[
\frac{x+7}{3}=9.
\]

\subsection*{Ваше решение}
В работе записано:
\[
x+7=9\cdot3,
\]
\[
x+7=27,
\]
\[
x=27-7,
\]
\[
x=20.
\]
После получения корня проверка подстановкой не выполнена.

\subsection*{Ваш ответ}
\[
x=20.
\]

\subsection*{Ошибка}
\ErrorLabel\ обязательная часть задания не выполнена: после нахождения корня отсутствует проверка.

\subsection*{Где именно возникает ошибка}
Все алгебраические преобразования выполнены верно, и корень $x=20$ найден правильно. Неполнота возникает после последней строки: условие требует проверить найденное значение в исходном уравнении.

\subsection*{Почему это неверно}
Если условие прямо требует проверку, решение нельзя считать завершённым сразу после нахождения корня. Нужно подставить найденное значение именно в исходное уравнение и убедиться, что левая и правая части равны.

Для $x=20$ получаем
\[
\frac{20+7}{3}=\frac{27}{3}=9.
\]
Левая часть равна правой, поэтому корень найден верно.

\subsection*{Ключевая идея}
\KeyIdea\ после нахождения корня подставьте его в исходное уравнение. Совпадение левой и правой частей завершает требуемую проверку.

\subsection*{Исправление от последнего правильного шага}
Последний правильный шаг уже есть:
\[
x=20.
\]
Добавляем недостающую проверку:
\[
\frac{20+7}{3}=\frac{27}{3}=9.
\]
Равенство выполняется.

\subsection*{Полное правильное решение}
Решаем уравнение:
\[
\frac{x+7}{3}=9.
\]
Умножаем обе части на $3$:
\[
x+7=27.
\]
Вычитаем $7$:
\[
x=20.
\]

Проверяем найденное значение в исходном уравнении:
\[
\frac{20+7}{3}=\frac{27}{3}=9.
\]
Левая часть равна правой.

\subsection*{Проверка}
\[
\frac{20+7}{3}=9.
\]
Равенство выполняется.

\subsection*{Итог}
\Outcome\ корень $x=20$ найден правильно. Для полного выполнения задания нужно было дописать проверку подстановкой.

\subsection*{Задача на закрепление}
Решите уравнение и выполните проверку:
\[
\frac{x-5}{4}=7.
\]

\newpage
\phantomsection\label{task:5b}
\section*{Задание 5б}

\subsection*{Текст задания}
Решите уравнение. Для пункта б сначала укажите ограничение на значение $x$, затем выполните проверку:
\[
\frac{12}{x-2}=3.
\]

\subsection*{Ваше решение}
В работе записано:
\[
\begin{aligned}
x-2&=12:3,\\
x-2&=4,\\
x&=4+2,\\
x&=6.
\end{aligned}
\]
Ограничение на значение $x$ перед решением не указано. Проверка найденного корня также отсутствует.

\subsection*{Ваш ответ}
\[
x=6.
\]

\subsection*{Ошибка}
\ErrorLabel\ пропущены два обязательных элемента: ограничение на значение $x$ до решения и проверка после нахождения корня.

\subsection*{Где именно возникает ошибка}
Первый недостающий шаг должен появиться ещё до преобразования уравнения. В знаменателе стоит выражение $x-2$, а знаменатель не может быть равен нулю. Поэтому нужно записать
\[
x-2\ne0,
\qquad
x\ne2.
\]

После этого выполненные преобразования приводят к правильному корню $x=6$, но решение снова заканчивается раньше требуемого: подстановка в исходное уравнение не выполнена.

\subsection*{Почему это неверно}
При $x=2$ знаменатель $x-2$ равен нулю, поэтому исходная дробь не имеет смысла. Значение $2$ необходимо исключить до дальнейших преобразований.

После записи ограничения основной ход решения приводит к допустимому корню $x=6$, поскольку $6\ne2$. Но условие требует ещё и подстановку найденного значения в исходное уравнение; без неё решение остаётся неполным.

\subsection*{Ключевая идея}
\KeyIdea\ если переменная находится в знаменателе, сначала исключите значение, при котором знаменатель равен нулю. Затем решите уравнение и проверьте найденный корень в исходной записи.

\subsection*{Исправление от последнего правильного шага}
Перед решением добавляем ограничение, затем сохраняем основной ход:
\[
\begin{aligned}
x-2&\ne0, & x&\ne2,\\
x-2&=12:3=4, & x&=6.
\end{aligned}
\]
Корень удовлетворяет ограничению. После этого обязательно подставляем $x=6$ в исходное уравнение и получаем $3$ в левой части.

\subsection*{Полное правильное решение}
Сначала устанавливаем ограничение:
\[
x-2\ne0,
\qquad
x\ne2.
\]
Решаем уравнение. Так как $x-2\ne0$, умножаем обе части на $x-2$:
\[
\begin{aligned}
12&=3(x-2),\\
4&=x-2,\\
x&=6.
\end{aligned}
\]
Полученное значение не нарушает ограничение $x\ne2$.

Проверяем корень в исходном уравнении:
\[
\frac{12}{6-2}=\frac{12}{4}=3.
\]
Левая часть равна правой.

\subsection*{Проверка}
Значение $x=6$ допустимо, потому что $6\ne2$, и
\[
\frac{12}{6-2}=3.
\]
Равенство выполняется.

\subsection*{Итог}
\Outcome\ корень $x=6$ найден правильно. Для полного решения перед преобразованиями нужно было записать $x\ne2$, а после нахождения корня --- выполнить проверку.

\subsection*{Задача на закрепление}
Решите уравнение, сначала указав ограничение на значение $x$, и выполните проверку:
\[
\frac{18}{x+1}=6.
\]

\end{document}
